% Source: https://artofproblemsolving.com/community/c3h3479722_amc_8_problem
% Stable identity: aops-topic-3479722-q01
Area ratios are unchanged by an invertible affine coordinate change, so take $A=(0,0),C=(3,0),B=(0,3)$. Then $D=(1,0)$ and $E=(\dfrac{1}{2},\dfrac{3}{2})$. Line $AE$ has equation $y=3x$; intersecting $BC$, where $x+y=3$, gives $F=(\dfrac{3}{4},\dfrac{9}{4})$. Thus $BF:BC=1:4$. The distance of $E$ to $BC$ is one-third that of $A$, since the linear quantity $3-x-y$ is $1$ at $E$ and $3$ at $A$. Consequently $[EBF]=\dfrac{1}{4}\cdot\dfrac{1}{3}[ABC]=\boxed{30}$.
\begin{asy}
usepackage("amsmath");
size(7cm);
pair A=(0,0),B=(0,3),C=(3,0),D=(1,0),E=(.5,1.5),F=(.75,2.25);draw(A--B--C--cycle);draw(B--D);draw(A--F);fill(E--B--F--cycle,lightgray);draw(E--B--F--cycle);label("$A$",A,SW);label("$B$",B,N);label("$C$",C,SE);label("$D$",D,S);label("$E$",E,W);label("$F$",F,NE);
\end{asy}
